% This is the configuration file for letters processed with ucsdletter.cls
% The entries below are used to make the data for the default subhead.
% By the subhead we mean the block containing the return address, telephone,
% email and so on. You may use up to 3 columns, \colm{A}, \colm{B} and \colm{C}.
% Each \colm is made up of an arbitrary number of \itm's.
% The \itm's will be by default left aligned in the \colm. 
% To right align a \colm, add [r] after \colm.
% You do not need to add \<space> after a period in these \itm's.
%
% UCSD style gidelines call for subhead to be Times-Roman at 8pt.
% This is defined in the class file, but may be changed here by uncommenting
% the next two lines. 
%\def\subheadfontname{ptmr at 8pt}% 
%\font\subheadfont=\subheadfontname
%
% You should change the \subheadbaselineskip from its default 10pt
% if you change the size of the subheadfont. Two points larger than fontsize is % typical. UCSD style guidelines call for \subheadbaselineskip to be 10pt.
%
%\setlength\subheadbaselineskip{10pt}
%
\colm{A}{% UCSD STYLE GUIDELINES SAY ENTRIES HERE SHOULD BE CAPS
\itm{Assistant Professor Robert J. Webber}
\itm{Department of Mathematics}
\itm{9500 Gilman Drive \# 0112}
\itm{La Jolla, California \ 92093--0112}
}
\colm{B}{%
\itm{EMAIL:}
\itm{URL:}
%\itm{ICBM:}
}
\colm{C}{%
\itm{rwebber@ucsd.edu}
\itm{https://sites.google.com/ucsd.edu/rwebber/}
%\itm{32${}^{\circ}$ 52' 44'' N, 117${}^{\circ}$ 14' 28'' W}
}
% Signature
% Set your signature as it should appear underneath the letter \Closing line:
%
\signature{Robert J. Webber\\ Assistant Professor of Mathematics}
%
% By filling out the correct entries above, you will get a subhead conforming
% to the UCSD style guidelines, and you may ignore the settings below.
% If you wish to customize your configuration further, read on.
%
% Within the subhead, indentation from left and right are controlled by 
% \subheadlefti and \subheadrighti.
% Separation between \colm B and \colm C is \subheadcolBCsep
%\subheadlefti=0pt
%\subheadrighti=0pt
\subheadcolBCsep=.1in% use .1in if you have data in both \colB and \colC
% End of canned subhead
%
% The default textwidth of a letter is 6in, with equal margins on left and
% right.
% To change the texwidth of the body of your letter, uncomment and change the 
% following:
%
%\setlength\textwidth{5in}
%
% To use a seal other than the default seal file, named ucseal, enter name:
%
%\def\sealfile{ucsealB} 
%
% don't add the extension on the filename---you should have a pdf version for 
% use in pdftex, an eps version if your TeX generates a dvi intermediary.
%
% The top of the logo should print just over .5 inches from the top of the 
% paper so that the line above the campus names is exactly 1.0in from the
% top of the paper. If your printer sets the logo too low on the page, decrease
% this amount accordingly.
% (Increase it to lower the logo.)
%
\setlength\logotop{44pt}
%
% The default behavior of the standard letter class is to add glue above and
% below a short first page, not quite centering it vertically. If you wish to
% defeat the vertical centering, uncomment the following line:
%
\vcenteringfalse
%
% The default is to make the subhead a part of the letterhead. 
% Uncomment the following line to change this, making it print at top of text.
% It may look strange if you set \subheadinheadfalse without 
% \vcenteringfalse as well.
%
%\subheadinheadfalse
%
% If the subhead is part of the letterhead, the default separation is .125in
% from the bottom of the seal to the top of the subhead.
% You can change this by uncommenting and completing the following:
%
%\setlength\subheadtopsep{.125in}
%
%
% Beginning of custom subhead
% If you wish to bypass the subhead described above in favor of a custom 
% subhead, uncomment the \customsubheadtrue line below.
%
%\customsubheadtrue
%
% You may add your own subhead, by adding the LaTeX code inside the following
% macro. Whether or not you wish it to print as part of the head or part of the
% text, the content of \mysubhead should be a \vbox with your data spelled out
% exactly. If you uncommented \subheadinheadfalse, the \vbox will be inserted
% at the top of the body of the letter. The code you enter will be wrapped in an
% \hbox of appropriate size and inserted as instructed.
% You will have to specify the \baselineskip, as the \baselineskip is set to 0
% when this box is constructed..
\def\mysubhead{}
%
% End of custom subhead
%
% Additional packages and settings
%
% (a) Within each address environment, such as in the subhead, the recipient
% address, signature, cc, encl and labels, the class sets \sfcode`.=1000, so 
% there is no need to modify the space following a period. For example, you may
% write Dr. Hancock rather than Dr.\ Hancock within those environments.
%
% (b) You may add here any additional packages and settings you wish.
% The effect is the same as adding them just after the 
% \documentclass{ucsdletter} line. They will execute before other options chosen
% in your letter, and may be overwritten by those later settings.
% You could, for example, select the font to be used in the body of the letter
% by using an appropriate \usepackage{}, or choose the form of the headings and % footers by a \pagestyle entry. (The default is plain, the other options being
% headings and empty. 
\AtEndOfClass{
\usepackage{mathptmx}% Times Roman text, Symbol math (may use Nimbus clones)
%\pagestyle{headings}%UCSD style guidelines call for this:
% other options are plain, empty
}
